\section{Zadaci}

\newtheorem{zadatak}{Zadatak}

\begin{zadatak}
    Odrediti napon $V_{out}$ za kolo sa slike \ref{fig:zad1}, ako su $R_1 = 4 \,k\,\Omega$ i $R_2 = 6\,k\,\Omega$, a
    \begin{enumerate}
        \item [a)] $V_{in} = 10 \,V$
        \item [b)] $V_{in} = 10\sin{(20\pi\,t)} \,V$
    \end{enumerate}
\end{zadatak}

\begin{figure}[ht]
\centering
\begin{circuitikz}[american]
    \draw (0,0) node[ground]{} to [V, invert,l=$V_{in}$] (0,3) to[R, l^=$R_1$, v_>=$V_{R_2}$] (6,3) to[R, l^=$R_2$, v_>=$V_{out}$] (6,0) node[ground]{};
    \draw[->] (0,3) -- node[above] {$I_{in}$} (1,3);
\end{circuitikz}
\caption{Kolo za zadatak 1}
\label{fig:zad1}
\end{figure}

\begin{zadatak}
    Odrediti vrednost za izlazni napon $V_{out}$ za kolo sa slike \ref{fig:zad2}.\newline
    \textit{Napomena: Brojevi nisu dati, jer se očekuje opšti izraz za $V_{out}$ ovog kola.}
\end{zadatak}

\begin{figure}[ht]
\centering
\begin{circuitikz}[american]
    \draw (0,0) node[ground]{} to [V, invert,l=$V_{in}$] (0,3) to[R, l^=$R_1$] (5,3) to[R, l^=$R_3$] (10,3) to[R, l^=$R_4$, v_>=$V_{out}$] (10,0) node[ground]{};
    \draw (5,3) to[R, l^=$R_2$] (5,0) node[ground]{};
\end{circuitikz}
\caption{Kolo za zadatak 2}
\label{fig:zad2}
\end{figure}

\pagebreak
\section{Zadaci za vežbu}

\begin{zadatak}
    Za naponski razdelnik sa slike \ref{fig:zad3} poznato je da je
    $V_{\mathrm{in}} = 12\,\mathrm{V}$. Potrebno je odabrati otpornosti
    tako da izlazni napon bude $V_{\mathrm{out}} = 3\,\mathrm{V}$,
    a struja kroz neopterećeni razdelnik
    $I_{\mathrm{in}} = 0{,}5\,\mathrm{mA}$.

    \begin{enumerate}
        \item[a)] Odrediti odnos otpornosti $R_1/R_2$.
        \item[b)] Odrediti vrednosti otpornosti $R_1$ i $R_2$.
        \item[c)] Odrediti izlazni napon $V_{out}$ ako otpornici
        $R_1$ i $R_2$ zamene mesta.
    \end{enumerate}
\end{zadatak}

\begin{figure}[H]
    \centering
    \begin{circuitikz}[american]
        \draw
            (0,0) node[ground] {}
            to[V, invert,l=$V_{in}$] (0,2)
            to[R, l=$R_1$, my voltage] (5,2)
            to[R, l=$R_2$, v=$V_{out}$, my voltage] (5,0)
            (5,0) node[ground] {};

    \end{circuitikz}
    \caption{Kolo za zadatak 3}
    \label{fig:zad3}
\end{figure}

\begin{zadatak}
    Na izlaznim priključcima $a$ i $b$ nepoznatog linearnog
    kola izvršena su dva merenja idealnim voltmetrom.
    Svi naponi posmatraju se od priključka $a$ prema
    priključku $b$, odnosno $V_{ab} = V_a - V_b$.

    Kada na izlaz nije priključeno opterećenje, izmeren je
    napon praznog hoda:
    \[
        V_{\mathrm{oc}} = 6\,\mathrm{V}.
    \]
    Nakon priključivanja otpornika
    $R_L = 3\,k\Omega$ između priključaka $a$ i $b$,
    izmeren je napon na opterećenju:
    \[
        V_L = 4\,\mathrm{V}.
    \]

    Na osnovu rezultata merenja:
    \begin{enumerate}
        \item[a)] Odrediti parametre Tevenenovog ekvivalenta
        $E_{\mathrm{th}}$ i $R_{\mathrm{th}}$.

        \item[b)] Odrediti parametre Nortonovog ekvivalenta
        $I_N$ i $R_N$.

        \item[c)] Odrediti napon i struju opterećenja
        kada se postojeći otpornik zameni otpornikom
        otpornosti $1{,}5\,k\Omega$.
        Pozitivan smer struje kroz opterećenje usvojiti
        od priključka $a$ ka priključku $b$.
    \end{enumerate}
\end{zadatak}

\pagebreak
\section{Zadatak za domaći}

\begin{zadatak}
    Odrediti vrednost izlaznog napona $V_{out}$ za kolo sa slike \ref{fig:zadd},
    ako su:
    \[
        V_1 = 12\,V\,,
        \qquad V_2 = 8\,V\,,
    \]
    \[
        R_1 = R_5 = 3\,k\Omega\,,
        \qquad R_2 = R_3 = R_4 = 5\,k\Omega.
    \]
\end{zadatak}

\begin{figure}[ht]
\centering
\begin{circuitikz}[american]
    \draw (0,0) node[ground]{} to [V, invert,l=$V_{1}$] (0,3) to[R, l^=$R_1$] (5,3) to[R, l^=$R_3$] (10,3) to[R, l^=$R_4$, v_>=$V_{out}$] (10,0) node[ground]{};
    \draw (5,3) to[R, l^=$R_2$] (5,0) node[ground]{};
    \draw (0,3) to[V, invert, l=$V_2$] (0,6) -- (5,6) to[R, l=$R_5$] (5,3);
\end{circuitikz}
\caption{Kolo za domaći}
\label{fig:zadd}
\end{figure}